\documentclass[11pt,a4paper]{article}

\usepackage[T1]{fontenc}
\usepackage[utf8]{inputenc}

% Stabilize PDF metadata in addition to SOURCE_DATE_EPOCH from build.py.
\ifdefined\pdfinfoomitdate\pdfinfoomitdate=1\fi
\ifdefined\pdftrailerid\pdftrailerid{}\fi
\ifdefined\pdfsuppressptexinfo\pdfsuppressptexinfo=-1\fi

\setlength{\oddsidemargin}{0mm}
\setlength{\evensidemargin}{0mm}
\setlength{\textwidth}{159mm}
\setlength{\topmargin}{-10mm}
\setlength{\textheight}{235mm}
\setlength{\parindent}{0pt}
\setlength{\parskip}{0.7em}
\raggedbottom
\sloppy

\newcommand{\code}[1]{\texttt{#1}}
\newcommand{\PipelineObserve}{Observe\\project}
\newcommand{\PipelinePlan}{Plan\\difference}
\newcommand{\PipelineStage}{Stage in\\isolation}
\newcommand{\PipelineVerify}{Verify\\candidate}
\newcommand{\PipelineFinish}{Commit or\\roll back}
\newcommand{\PipelineFailure}{Failure: return unchanged}
\newcommand{\ProjectBoundaryTitle}{Project-owned}
\newcommand{\ProjectBoundaryOne}{application sources}
\newcommand{\ProjectBoundaryTwo}{domain configuration}
\newcommand{\ProjectBoundaryThree}{unknown files}
\newcommand{\ProjectBoundaryFour}{product repositories}
\newcommand{\ToolingBoundaryTitle}{Tooling-owned}
\newcommand{\ToolingBoundaryOne}{profiles and rules}
\newcommand{\ToolingBoundaryTwo}{adapters and resources}
\newcommand{\ToolingBoundaryThree}{migrations and state}
\newcommand{\ToolingBoundaryFour}{tooling documentation}
\newcommand{\BoundaryObserve}{evidence}
\newcommand{\BoundaryProposal}{bounded patch}
\newcommand{\BoundaryContract}{explicit ownership contract}

\title{Portable Tooling Architecture\\A Bilingual Case Study of Controlled Integration}
\author{Portable tooling case study}
\date{Reproducible edition}

\begin{document}
\maketitle

\begin{abstract}
This case study examines a tooling package that can be copied directly into different
software projects. Unlike a product template, it owns neither the application nor its
lifecycle. It observes existing technologies, selects a capability profile, plans bounded
changes and applies them as a transaction. The evaluation focuses on portability, ownership
protection, deterministic execution and the cost of a deliberately fail-closed design. The
results suggest that the strongest guarantees come not from a large collection of generators,
but from explicit boundaries, repeatable plans and demonstrable recovery.
\end{abstract}

\tableofcontents
\newpage

\section{Context and research question}

Many internal tooling repositories begin as complete project templates. Over time they
accumulate assumptions about the frontend, backend, database, desktop shell, CI and release
process. This is convenient while every target follows the same layout. Once the tooling is
copied into an existing or differently structured project, those assumptions become risks:
fixed paths point at the wrong files, a generator overwrites product code, and an upgrade can
no longer distinguish its own state from state owned by the application team.

The architecture studied here replaces that model with a portable tooling directory. Its
central research question is:

\begin{quote}
How can copyable tooling support several project shapes without owning application sources,
while making integration, migration and verification reproducible and reversible?
\end{quote}

The question is broader than packaging. Portability means that paths originate in a project
context resolved at runtime, technology knowledge resides behind adapters, and every write is
visible in advance as a structured plan. A copied directory is only truly portable if its
behavior is independent of the location and layout of the repository in which it was
developed.

\section{Architecture of the portable core}

\subsection{Context instead of a global root}

The project context separates the project root, tooling root, documentation root, packaged
resources and local state. Default roots are derived centrally from the installed tooling
location, validated at the boundary and then passed explicitly. Consumers therefore do not
each guess a separate global repository root. This distinction is small in code but fundamental
in operation: the same copy may live inside a monorepo, next to a web application or in a
desktop project without recompiling source paths.

Project configuration describes paths and profile intent but remains project-owned. Profiles,
adapter contracts and canonical resources belong to the tooling. Local integration state is
kept outside the tooling directory in a bounded state area. As a result, the complete tooling
directory can be replaced without mixing project history with an old implementation.

\subsection{Profiles and adapters}

A profile expresses desired capabilities rather than a complete directory tree. Examples
include a web-only setup, a cloud-enabled variant and a desktop combination. Adapters map
those capabilities to technology-specific observations and controlled actions. A frontend
adapter may recognize an existing package manifest, while a database adapter activates only
when strong markers are present. With insufficient evidence an adapter remains inactive;
speculation never grants write access.

This separation keeps the core narrow. It understands states, findings, plans, transactions
and migrations. It does not treat particular npm scripts or database paths as universal
truths. A new technology extends the adapter registry without forcing the core into a new
product model.

\section{Ownership as a safety boundary}

Figure~\ref{fig:ownership} shows the architecture's most important boundary. The project owns
its application, domain configuration and every unknown file. The tooling owns only its
packaged rules, resources, migrations and documentation. An adapter may observe project files
and propose bounded changes; observation does not imply blanket ownership of the target tree.

\begin{figure}[htbp]
  \centering
  % Portable, package-free LaTeX picture. Labels are supplied by each language source.
\begingroup
\setlength{\unitlength}{1mm}
\begin{picture}(118,54)
  \put(0,7){\framebox(50,40){%
    \parbox{43mm}{\centering\small\textbf{\ProjectBoundaryTitle}\\[2mm]
      \scriptsize\ProjectBoundaryOne\\\ProjectBoundaryTwo\\
      \ProjectBoundaryThree\\\ProjectBoundaryFour}}}
  \put(68,7){\framebox(50,40){%
    \parbox{43mm}{\centering\small\textbf{\ToolingBoundaryTitle}\\[2mm]
      \scriptsize\ToolingBoundaryOne\\\ToolingBoundaryTwo\\
      \ToolingBoundaryThree\\\ToolingBoundaryFour}}}
  \put(50,31){\vector(1,0){18}}
  \put(59,35){\makebox(0,0){\scriptsize\BoundaryObserve}}
  \put(68,20){\vector(-1,0){18}}
  \put(59,16){\makebox(0,0){\scriptsize\BoundaryProposal}}
  \put(59,51){\makebox(0,0){\small\textbf{\BoundaryContract}}}
\end{picture}
\endgroup

  \caption{Evidence flows to the tooling; only bounded proposals flow back.}
  \label{fig:ownership}
\end{figure}

The ownership contract operates at three levels:

\begin{enumerate}
  \item \textbf{Path boundaries:} relative paths are normalized, must remain below the
  project root and may not exploit symbolic indirection.
  \item \textbf{Change boundaries:} structured operations distinguish creation,
  replacement and removal. Preconditions stop a changed file from being silently
  overwritten between observation and commit.
  \item \textbf{Knowledge boundaries:} adapters declare their markers and target paths.
  Unknown files are not errors and are not treated as debris to be cleaned up.
\end{enumerate}

These rules are stricter than a conventional generator. They reduce the scope of automatic
repair while increasing traceability. For a tool that runs in repositories it does not own,
that asymmetry is desirable: a skipped convenience step can be corrected, whereas lost
product work may not be recoverable.

\section{Transactional integration}

The pipeline in Figure~\ref{fig:pipeline} separates observation from publication. Current
state, profile intent and adapter findings first produce a deterministic plan. The planned
operations are then applied to an isolated staging copy. When the plan requires them, fixed
tooling quality, test and locked-dependency checks execute against that candidate. Only a
completely successful candidate may be published to the project.

\begin{figure}[htbp]
  \centering
  % Portable, package-free LaTeX picture. Labels are supplied by each language source.
\begingroup
\setlength{\unitlength}{1mm}
\begin{picture}(118,31)
  \put(0,12){\framebox(20,12){\scriptsize\shortstack{\PipelineObserve}}}
  \put(20,18){\vector(1,0){4}}
  \put(24,12){\framebox(20,12){\scriptsize\shortstack{\PipelinePlan}}}
  \put(44,18){\vector(1,0){4}}
  \put(48,12){\framebox(20,12){\scriptsize\shortstack{\PipelineStage}}}
  \put(68,18){\vector(1,0){4}}
  \put(72,12){\framebox(20,12){\scriptsize\shortstack{\PipelineVerify}}}
  \put(92,18){\vector(1,0){4}}
  \put(96,12){\framebox(22,12){\scriptsize\shortstack{\PipelineFinish}}}
  \put(83,10){\vector(0,-1){5}}
  \put(83,2){\makebox(0,0){\scriptsize\PipelineFailure}}
  \put(83,5){\line(-1,0){59}}
  \put(24,5){\vector(0,1){7}}
\end{picture}
\endgroup

  \caption{Integration publishes only a fully verified candidate.}
  \label{fig:pipeline}
\end{figure}

The commit is controlled as well. Preconditions observed during planning are checked again
before writes begin. Backups and a transaction journal support rollback if a later operation
fails. Persisted state does not store a blanket snapshot of the project; it stores only the
evidence needed for ownership and drift decisions. A file edited by a maintainer therefore
remains visible instead of being replaced as if it were an outdated template copy.

Two modes follow from the same planning model. A read-only check produces findings and a plan
without changing a byte. A full repair may execute only operations authorized by that plan.
After a successful run, another check must report no drift, and a second repair must be a
no-op. This idempotence requirement is stronger evidence than the one-time exit code of a
generator.

\section{Executing external toolchains}

Some verification requires real toolchains. A staged candidate can require trusted tooling
quality and test checks or locked-dependency validation through offline
\code{npm ci --ignore-scripts}. Separately, a user can explicitly dispatch a fixed live
adapter capability for an existing product. The tooling cannot replace the semantics of
those programs, but it can bound their invocation; arbitrary user-configured shell fragments
are outside the contract.

Both action paths run without a shell, with bounded output, a timeout and disposable home,
cache and configuration directories. Staged validation additionally forces supported
toolchains offline and disables lifecycle scripts; an explicit live action may intentionally
contact registries and create project artifacts. Process groups or Job Objects terminate
descendants when an action ends. These controls are not an operating-system sandbox: executed
code still has the local user's file permissions. They make the remaining trust assumption
explicit, while only staged validation belongs to the integration rollback boundary.

\section{Evaluation}

\subsection{Experimental design}

The architecture is evaluated with five profiles in independent temporary target projects.
Each profile follows the same sequence: copy the tooling, detect project state, check, perform
the complete integration, check again, execute applicable tests and run integration a second
time. A complementary experiment integrates the pinned historical 0.1.0 payload and then
replaces both managed directories with the current release to exercise a registered migration
and the separation between implementation and project state.

The measures are deliberately binary and byte-oriented:

\begin{itemize}
  \item Hashes of all unknown and product-owned files remain unchanged.
  \item A read-only check changes neither files nor directory structure.
  \item After successful integration, the second integration plans zero operations.
  \item An intentionally failing verification restores the initial state.
  \item The same portable payload retains its manifest identity after a direct copy.
\end{itemize}

\subsection{Results and interpretation}

\begin{center}
\begin{tabular}{|p{43mm}|p{48mm}|p{55mm}|}
\hline
\textbf{Criterion} & \textbf{Observable evidence} & \textbf{Interpretation} \\
\hline
Relocation & Context paths point only into the temporary target & No dependency on the source
repository \\
\hline
Ownership & Product and unknown-file hashes remain stable & Adapter knowledge bounds rather
than expands ownership \\
\hline
Idempotence & A second complete run plans no operations & State and desired configuration
agree \\
\hline
Rollback & A failed candidate leaves no partial change & Verification precedes publication \\
\hline
Tooling replacement & Migration uses preserved project state & Implementation and state are
separate \\
\hline
\end{tabular}
\end{center}

The evaluation does not prove that every conceivable adapter is correct. It demonstrates the
properties that every adapter must preserve. The transaction core can roll back execution or
verification failures, but it cannot prove that an otherwise valid allowlisted domain change
is semantically desirable. Conversely, a domain-correct adapter can break portability by using
fixed root paths. The test strategy therefore combines profile matrices with contract tests
for paths, markers, ownership scopes and action declarations.

\section{Trade-offs and limitations}

\textbf{More preparation, smaller failure surface.} An adapter needs declared markers,
operations and tests. This costs more than copying a prepared directory tree, but it avoids
implicit ownership claims.

\textbf{Fail-closed behavior may reduce availability.} If a lockfile, compiler or safe path
resolution is missing, the tooling refuses the automatic change. Teams must resolve such
findings deliberately. The benefit is that uncertainty is never interpreted as permission.

\textbf{Reproducibility remains environment-sensitive.} Fixed time, locale, offline flags and
clean caches remove many sources of drift. Different compiler or toolchain versions can still
produce different artifacts. Full byte reproducibility also requires a versioned build
environment.

\textbf{Process containment is not a sandbox.} Timeouts and process trees stop common leaks,
but they do not isolate filesystem permissions from malicious tests. Untrusted projects need
an external virtual-machine or container boundary.

\textbf{The manifest is a self-consistency boundary.} An exact payload inventory detects mixed
or changed tooling copies only while the bundled manifest remains trusted. Coordinated changes
to payload and manifest are not authenticated. It must be regenerated after every legitimate
payload change; a stale manifest blocks portable integration and verification, while a release
gate must invoke that validation explicitly.

\section{Conclusion}

Moving from a master template to portable tooling is primarily a change in ownership. Profiles
describe capabilities, adapters provide bounded technology knowledge, and the core publishes
only verified transactions. Explicit context creates portability; narrow paths and fail-closed
decisions create safety; separating project state from the replaceable implementation creates
maintainability.

The case study also marks the limits of that claim. The system cannot make arbitrary external
code trustworthy, nor can it equalize toolchain versions without a controlled build
environment. It can make each assumption visible, prove the scope of every change and return
to a known state after failure. For tooling intended to operate across many independently
evolved projects, that controlled restraint is the central architectural achievement.

\end{document}
